\documentclass[11pt,a4paper]{article}
% Kodowanie i obsługa języka polskiego
\usepackage[utf8]{inputenc}      % Kodowanie wejścia
\usepackage[T1]{fontenc}         % Kodowanie fontów
\usepackage[polish]{babel}       % Obsługa języka polskiego

% Pakiety matematyczne i inne przydatne
\usepackage{mathtools}
\usepackage{tabularx}
\usepackage{array} % for centering in p-columns
\usepackage{calc}
\usepackage[makeroom]{cancel}
\usepackage{xcolor}
\usepackage{bm}
\usepackage{marginnote}
\usepackage{tikz}
\usetikzlibrary{trees}
\usetikzlibrary{decorations.pathreplacing}
\usepackage{amsmath, amssymb, amsthm}  % Pakiety do matematyki
\usepackage{graphicx}                  % Obsługa grafiki
\usepackage{hyperref}                  % Linki i spis treści
\usepackage{geometry}                  % Ustawienie marginesów
\usepackage{algorithm}               % Algorytmy
\usepackage{algpseudocode}           % Pseudokod
\newcommand{\curly}{\mathrel{\leadsto}}

% Twierdzenia, lematy itd.
\newtheorem{twierdzenie}{Twierdzenie}[section]
\newtheorem{fakt}[twierdzenie]{Fakt}
\newtheorem{lemat}[twierdzenie]{Lemat}
\newtheorem{wniosek}[twierdzenie]{Wniosek}
\newtheorem{stwierdzenie}[twierdzenie]{Stwierdzenie}

% Definicje, przykłady, uwagi
\theoremstyle{definition}
\newtheorem{definicja}[twierdzenie]{Definicja}
\newtheorem{przyklad}[twierdzenie]{Przykład}
\newtheorem{zadanie}[twierdzenie]{Zadanie}

\theoremstyle{remark}
\newtheorem{uwaga}[twierdzenie]{Uwaga}
\newtheorem{oznaczenie}[twierdzenie]{Oznaczenie}
% Dodatkowe ustawienia
\newlength{\exprwidth}
\hypersetup{
    colorlinks=true,
    linkcolor=blue,
    filecolor=magenta,
    urlcolor=cyan,
}
\title{Notatki z Wybranych zagadnień Algebry}
\author{Jakub Kogut}

\begin{document}

\maketitle

\tableofcontents  % Spis treści (opcjonalnie)
\newpage

\section{Wstęp}
Notatki z przedmiotu Wybrane zagadnienia Algebry. Prowadzony przez dr Bojko.
\section{Wyklad 04.03.2026}
Krótkie przypomnenie algebry z poprzednich lat.
\subsection{Podstawowe definicje}
\begin{definicja}[Ciało, Field]~\\
    Jeśli Pierścień z jedynka spełnia następujące warunki:
    \begin{itemize}
        \item $(X\setminus \{0\}, \cdot)$ jest grupą abelową
    \end{itemize}
    to nazywamy go ciałem.
\end{definicja}
\begin{definicja}[Rozszerzenie pierścienia]~\\
    Mamy pierścień $\mathcal{R} = (X, +, \cdot, 0)$ oraz $A$ taki, że $A \cap X = \emptyset$. Rozszerzeniem pierścienia $\mathcal{R}$ o $A$ nazywamy pierścień $\mathcal{R}' = (X', +', \cdot', 0')$, gdzie $X' = X \cup A$, a operacje $+'$ i $*'$ są rozszerzeniami odpowiednio $+$ i $\cdot$.
\end{definicja}
\begin{uwaga}~\\
    Jeśli $|A| = n \in \mathbb{N}$, to dla $A= \{a_1, \dots, a_n\}$, piszemy 
    \[
        \mathcal{R} [A]=\mathcal{R}[a_1, \dots, a_n]
    \]
\end{uwaga}
\begin{przyklad}[Pierścień Gaussa]~\\
    $\mathbb{Z}[i]$ jest rozszerzeniem pierścienia $\mathbb{Z}$ o element $i$, gdzie $i^2 = -1$. W tym przypadku, $\mathbb{Z}[i]$ składa się z wszystkich liczb postaci $a + bi$, gdzie $a, b \in \mathbb{Z}$.
\end{przyklad}
\begin{przyklad}~\\
    $\mathbb{Q}[\sqrt{2}]$ jest rozszerzeniem pierścienia $\mathbb{Q}$ o element $\sqrt{2}$, gdzie $(\sqrt{2})^2 = 2$. W tym przypadku, $\mathbb{Q}[\sqrt{2}]$ składa się z wszystkich liczb postaci $a + b\sqrt{2}$, gdzie $a, b \in \mathbb{Q}$.
\end{przyklad}
\begin{przyklad}[Liczby Dualne]~\\
    $(D, +, \cdot, 0, 1)$, gdzie $D = \{ a + b\epsilon : a, b \in \mathbb{R}, \epsilon^2 = 0 \}$, jest rozszerzeniem pierścienia $\mathbb{R}$ o element $\epsilon$, który jest nilpotentny (tzn. $\epsilon^2 = 0$). W tym przypadku, $D$ składa się z wszystkich liczb postaci $a + b\epsilon$, gdzie $a, b \in \mathbb{R}$.
\end{przyklad}
\begin{przyklad}[Liczby Podwójne]~\\
$(D, +, \cdot, 0, 1)$, gdzie $D = \{ a + bj : a, b \in \mathbb{R}, j^2 = 1 \}$, jest rozszerzeniem pierścienia $\mathbb{R}$ o element $j$, który jest idempotentny (tzn. $j^2 = 1$). W tym przypadku, $D$ składa się z wszystkich liczb postaci $a + bj$, gdzie $a, b \in \mathbb{R}$.
\end{przyklad}
\begin{definicja}[Dziedzina Euklidesowa]~\\
    Dziedziną całkowitości $\mathbb{R} = (X, +, \cdot, 0)$ nazywamy dziedzinę euklidesową, jeśli
    \[
        (\exists N: X \rightarrow \mathbb{N}) \forall a, b \in X: (b \neq 0 \Rightarrow \exists q, r \in X: a = bq + r \land (r = 0 \lor N(r) < N(b)))
    \]
\end{definicja}
\begin{uwaga}~\\
    Jeżeli istneje $N$, które spełnia warunek z definicji dziedziny euklidesowej, to można przyjąć, że $N$ jest minimalne, tzn. dla każdego $a \in X$ i każdego $N': X \rightarrow \mathbb{N}$, jeśli $N'$ spełnia warunek z definicji, to $N(a) \leq N'(a)$.
\end{uwaga}

\section{Wykład 11.03.2026}
Na start krótka poprawka do poprzedniego wykładu.
\subsection{Dziedziny całkowitości}
\begin{definicja}~\\
Dziedziną całkowitości $\mathbb{R} = (X, +, \cdot, 0)$ nazywamy dziedzinę euklidesową, jeśli
\[
    (\exists N: X \rightarrow \mathbb{N}) \forall a, b \in X: (b \neq 0 \Rightarrow \exists q, r \in X: a = bq + r \land (r = 0 \lor N(r) < N(b)))
\]
oraz
\[
    (\forall a, b \in X) (N(ab) \geq N(a))
\]
\end{definicja}
\begin{przyklad}~\\
    Dziedziną całkowitości $\mathbb{Z}$ jest $\mathbb{Z}$ z funkcją $N(a) = |a|$.
\end{przyklad}
\begin{przyklad}~\\
    Dziedziną całkowitości $\mathbb{Z}[i]$ jest $\mathbb{Z}[i]$ z funkcją $N(a + bi) = a^2 + b^2$.
\end{przyklad}
Rozwazmy izomorfizm 
\[
    \mathbb{C} \cong \left\{\begin{bmatrix}a & -b \\ b & a
\end{bmatrix} : a, b \in \mathbb{R}\right\}
        \]
        \[
    \varphi(a + bi) = \begin{bmatrix}a & -b \\ b & a
\end{bmatrix}
\]
można to udowodnić jako ćwiczenie
\[
    \varphi((a + bi)(c + di)) = \begin{bmatrix}ac - bd & -(ad + bc) \\ ad + bc & ac - bd\end{bmatrix} = \varphi(a + bi)\varphi(c + di)
\]
\begin{przyklad}~\\
    wracając do pierścienia Gaussa, zastanówmy się nad ideałem $<a + bi>$. Czy jest on główny? Odpowiedź brzmi tak, ponieważ $<a + bi> = <\gcd(a, b)>$.
\end{przyklad}
A jak wyglądają elementy odwracalne w tym pierścieniu? Są to elementy o normie równej 1, czyli $1, -1, i, -i$. Ponieważ $N(a + bi) = a^2 + b^2$, to $a^2 + b^2 = 1$ ma tylko te cztery rozwiązania całkowite.
\begin{przyklad}~\\
    Niech $K[x]$ będzie pierścieniem wielomianów nad ciałem $K$. Zdefiniujmy funkcję $N: K[x] \rightarrow \mathbb{N}$ jako $N(f) = \deg(f)$. Wtedy $K[x]$ jest dziedziną całkowitości, ponieważ dla dowolnych $f, g \in K[x]$ z $g \neq 0$, istnieją wielomiany $q, r \in K[x]$ takie, że $f = gq + r$ oraz $\deg(r) < \deg(g)$ lub $r = 0$. Ponadto, dla dowolnych $f, g \in K[x]$, mamy $\deg(fg) = \deg(f) + \deg(g) \geq \deg(f)$, co spełnia warunek drugiej części definicji dziedziny całkowitości.
\end{przyklad}
\begin{przyklad}~\\
    Rozważmy wielomian $w = x^{11}+10x \in \mathbb{Z}_{11}[x]$. Można łatwo załóważyć, ze jego zera to:
    \[
        w(0) = 0, w(1) = 1 + 10 = 0, \dots, w(10) = 10^{11} + 10 \cdot 10 = 10 + 100 = 0
    \]
    Wynika to z małego tw. Fermata.
\end{przyklad} 
Rodzi się pytanie jak sprawdzić, kiedy $f \in K[x]$ spełnia $f \equiv 0$?\\
Odpowiedź brzmi: $f \equiv 0$ wtedy i tylko wtedy, gdy $f$ ma co najmniej $|K|$ różnych zer. W przypadku $w$, ponieważ $w$ ma 11 różnych zer, to $w \equiv 0$ w $\mathbb{Z}_{11}[x]$.
\begin{fakt}~\\
    $f \equiv 0$ jeśli:
    \[
        \exists A \subseteq K: |A| > \deg(f) 
    \]
    oraz
    \[
    \forall a \in A: f(a) = 0
    \]
\end{fakt}
\begin{definicja}[Leading Term]~\\
    Niech $f \in K[x]$, $\sum_{i=0}^n a_i x^i$, $a_n \neq 0$. \\
    Określamy $LT(f) = a_n x^n$ jako \textit{najwyższy wyraz} wielomianu $f$, $LT(0) = 0$.
\end{definicja}
\subsection{Dziedziny ideałów głównych}
\begin{definicja}[Principal Ideal Domain, PID]~\\
    Dziedzine całkowitości $\mathcal{R}$ nazywamy dziedziną ideałów głównych, jeśli
    \[
        \forall \mathcal{I} \in ID(\mathcal{R}): \exists a \in \mathcal{R}: \mathcal{I} = \langle a \rangle
    \]
    gdzie $\langle a \rangle$ jest ideałem głowny.
\end{definicja}
\begin{fakt}
    W dziedzinach euklidesowych zachodzi:
    \[
        \langle a,b \rangle = \langle \gcd(a, b) \rangle
    \]
\end{fakt}
\begin{fakt}
    \[
        d|a \rightarrow \langle a \rangle \subseteq \langle d \rangle
    \]
\end{fakt}
\begin{wniosek}
    \[
        (d|a \land d|b) \rightarrow \langle a, b \rangle \subseteq \langle d \rangle
    \]
\end{wniosek}

\begin{fakt}
    W $\mathcal{R} \in EUCL$ ideały generowane przez $(a,b)$, $b \neq 0$, $a = bq + r$ wyglądają następująco:
    \[
        \langle a, b \rangle = \langle b, r \rangle
    \]
\end{fakt}
Aby to pokazać, należy zauważyć, że $ r \in \langle a, b \rangle$, ponieważ $r = a - bq$. Z drugiej strony, $a \in \langle b, r \rangle$, ponieważ $a = bq + r$. Ponadto, $b$ jest oczywiście elementem obu ideałów. W związku z tym, $\langle a, b \rangle$ i $\langle b, r \rangle$ zawierają te same elementy, co oznacza, że są one równe.
\subsection{Algorytm Euklidesa}
\begin{twierdzenie}[Euklidesa]
    Jeśli $NWD(n,m) = d$, to istnieją liczby całkowite $x, y$ takie, że $nx + my = d$. Jest to tożsamość Bézouta. Ponadto, $d$ jest najmniejszą dodatnią liczbą całkowitą, którą można przedstawić w postaci $nx + my$.
\end{twierdzenie}
\begin{algorithm}
\caption{Algorytm Euklidesa}
\begin{algorithmic}[1]
\Procedure{ExtendedGCD}{$a, b$}
\If{$b = 0$}
    \State \Return $(a, 1, 0)$
\Else
    \State $(d, x1, y1) \gets \text{ExtendedGCD}(b, a \mod b)$
    \State $x \gets y1$
    \State $y \gets x1 - \lfloor a / b \rfloor \cdot y1$
    \State \Return $(d, x, y)$
\EndIf
\EndProcedure
\end{algorithmic}
\end{algorithm}

\begin{przyklad}~\\
    Niech $f,g \in K[x]$,
    \begin{algorithmic}[1]
    \Procedure{div}{$f, g$}
    \State $q \gets 0$
    \State $p \gets f$
    \While{$p \neq 0 \&\& LT(g) | LT(p)$}
        \State $q \gets q + \frac{LT(p)}{LT(g)}$
        \State $p \gets p - \frac{LT(p)}{LT(g)}g$
    \EndWhile
    \State \Return $(q, p)$
    \EndProcedure
    \end{algorithmic}
    $q$ to wynik dzielenia, natomiast $p$ to reszta z dzielenia $f$ przez $g$. Algorytm ten jest poprawny, ponieważ w każdej iteracji zmniejszamy stopień wielomianu $p$, a proces kończy się, gdy $p$ jest równe zero lub gdy najwyższy wyraz $g$ nie dzieli najwyższego wyrazu $p$.
\end{przyklad}
\begin{definicja}[Relacja porządkowa]~\\
    W $\mathbb{N}^k$ definujemy relację porządkową $\leq$ w następujący sposób:
    \[
        (x_1, \dots, x_k) \leq (y_1, \dots, y_k) \iff \bigwedge_{i=0}^k: x_i \leq y_i
    \]
\end{definicja}
\begin{przyklad}~\\
    W $\mathbb{R}[x, y]$ 
    \[
        xy^2 \curly (1, 2) \leq x^2y^4 \curly (2, 4)
    \]
\end{przyklad}

\begin{lemat}[Dicksona]
    Niech $\overline{x_n} = \left( x_1^{(n)}, \dots, x_k^{(n)} \right)$ będzie nieskończoną sekwencją elementów z $\mathbb{N}^k$. Wtedy istnie $i_1, i_2, \dots$, że $(\overline{x_i})_{j \in \mathbb{N}}$ jest niemalejącą sekwencją.
\end{lemat}
\begin{proof}
    Niech $A_0 = \overline{x}[\mathbb{N}]$. Rozważmy następujące przypadki:
    \begin{enumerate}
        \item Jeśli $|A_0| < \aleph_0$, to istnie $\overline{z} \in \mathbb{N}^k$, że
        \[
            \left| \overline{z}^{-1}[\overline{z}] \right| = \aleph_0
        \]
        wystarczy wziąść podciąg taki, że $(\overline{x}_i_j) = \overline{z}$ dla każdego $j \in \mathbb{N}$, a wtedy $(\overline{x}_i_j) \in \mathbb{N}}$ jest niemalejącą sekwencją.
        \item Indukcja wzgledem $k$:
        \begin{itemize}
            \item Dla $k = 1$: \\
                kładziemy $\overline{x}_i_j = \min A_0$, $i_j \in \{ n \in \mathbb{N}: x_n = min A_0 \}$.
            \[
                A_1 = A_0 \setminus \{ \overline{x}_i_j \}
            \]
            Dalej $\overline{x}_i_j = \min A_{j-1}$, $i_j \in \{ n \in \mathbb{N}: x_n = min A_{j-1} \}$, $A_j = A_{j-1} \setminus \{ \overline{x}_i_j \}$.
            \[
                A_j = A_{j-1} \setminus \{ \overline{x}_i_j \}
            \]
            oczywiście $\overline{x}_i_j$ jest nierosnącą sekwencją.
        \item $(k \rightarrow k +1)$ \\
            Ciąg $(x_{k+1}^{(n)})_n$ ma podciąg niemalejący $(x_{k+1}^{(i_j)})_j$. \\
            Z założenia indukcyjnego
            \[
                y_j = ( x_1^{(i_j)}, \dots, x_k^{(i_j)} )
            \]
            istneje $(y_{j_l})_l$ niemalejący w $\mathbb{N}^k$. Zatem 
            \[
                \left( x_1^{(i_{j_l})}, \dots, x_k^{(i_{j_l})}, x_{k+1}^{(i_{j_l})} \right)
            \]
            jest niemalejący
    \end{itemize}
    \end{enumerate}
\end{proof}
\begin{wniosek}
    Jeśli $A \subseteq \mathbb{N}^k$, to zbiór elementów minimalnych w $A$ jest skończony.
\end{wniosek}
\section{Wykład 08.04.2026}
\begin{definicja}[Rozmaitości Algebraiczne]~\\
    Niech $\mathcal{F} \subseteq K[x_1,\dots,x_n]$ to zbiór 
    \[
        V(\mathcal{F}) = \{ (a_1, \dots, a_n)\in K^n: \forall f \in \mathcal{F}: f(\hat{a})=0\}
    \]
\end{definicja}
\begin{przyklad}
    \begin{enumerate}
        \item $V(\{0\}) = K^n$
        \item $ V(\{x, x+1\})$
    \end{enumerate}
\end{przyklad}
\begin{fakt}
    \begin{enumerate}
        \item $V(\mathcal{F}) \cup V(\mathcal{G}) = V(\{fg: f \in \mathcal{F}, g \in \mathcal{G}\})$
        \item $V(\mathcal{F}) \cap V(\mathcal{G}) = V(\mathcal{F} \cup \mathcal{G})$
    \end{enumerate}
\end{fakt}
\begin{przyklad}~\\
    $V(\{x^2 + y^2 - 1\})$ to okrąg o promieniu 1.
\end{przyklad}

\begin{definicja}[Przestrzeń Afiniczna]~\\
    Niech $A \subseteq K^n$, wtedy mówimy, że, 
    \[
        \mathcal{I}(A) = \{ f \in K[x_1, \dots, x_n]: \forall (a_1, \dots, a_n) \in A: f(a_1, \dots, a_n) = 0\}.
    \]
\end{definicja}
\begin{fakt}
    \begin{enumerate}
        \item $\mathcal{I}(A)$ to ideał
        \item Niech $I$ -- ideał. Wtedy $I \subseteq \mathcal{I}(V(I))$
    \end{enumerate}
\end{fakt}
\begin{przyklad}~\\
    Rozważmy ideał generowany przez
    \[
        \langle x^2, y^2 \rangle = \{a(x,y)x^2+b(x,y)y^2: a,b \in K[x,y]\}
    \]
    Wtedy $V(\lange x^2, y^2 \rangle) = \{(0, 0)\}$, a 
    \[
        \mathcal{I}(\{(0, 0)\}) = \langle x, y \rangle
    \]
    wszyskie kombinacje liniowe.
\end{przyklad}
\begin{definicja}[Porządek jedomianowy (MO)]~\\
    Dobry porządek na $\leq$ na $\mathbb{N}^n$ nazywamy porządkiem jedomianowym, jeśli 
    \[
        \forall \alpha, \beta, \gamma \in \mathbb{N}^n: (\alpha \leq \beta \Rightarrow \alpha + \gamma \leq \beta + \gamma)
    \]
    gdzie $a < b \iff a \leq b \land a \neq b$.
\end{definicja}
\begin{przyklad}[leksykograficzny]~\\
    $\alpha = (a_1, \dots, a_n) <_{lex} \beta = (b_1, \dots, b_n)$ wtedy i tylko wtedy, gdy istnieje $i \in \{1, \dots, n\}$ takie, że $a_j = b_j$ dla każdego $j < i$ oraz $a_i < b_i$.
\end{przyklad}
\begin{definicja}~\\
    Przy pomovy porządku jednomianowego MO na $\mathbb{N}^n$, możemy zbudować dobry porządek na jedenomianach $K[x_1, \dots, x_n]$ definiując 
    \[
        \bar{x}^\alpha \leq \bar{x}^\beta \iff \alpha \leq \beta
    \]
\end{definicja}
\begin{przyklad}~\\
    Dla $\leq_{lex}$ na $\mathbb{N}^2$, 
    \[
        1 < y < y^2 < \dots < x < xy < xy^2 < \dots < x^2 < x^2y < \dots
    \]
\end{przyklad}
\begin{definicja}[Graded Lexicographic Order]~\\
    $\alpha = (a_1, \dots, a_n) <_{grlex} \beta = (b_1, \dots, b_n)$ wtedy i tylko wtedy, gdy $|\alpha| < |\beta|$ lub $|\alpha| = |\beta|$ i $\alpha <_{lex} \beta$, gdzie $|\alpha| = a_1 + \dots + a_n$.
\end{definicja}
\begin{twierdzenie}~\\
    Jeśli $\leq$--LO na $\mathbb{N}^n$ oraz 
    \[
        \forall \alpha, \beta \in \mathbb{N}^n: \alpha < \beta \rightarrow \alpha + \gamma < \beta + \gamma
    \]
    to $\leq$--WO jest równowazne z warunkiem $\forall \alpha, \beta \in \mathbb{N}^n: \alpha < \beta \rightarrow \alpha + \gamma < \beta + \gamma$.
\end{twierdzenie}
\begin{proof}
    Załóżmy, że $\alpha <0$. Wtedy $\alpha + \alpha < \alpha + 0$. Analogicznie $(k+1)\alpha < k\alpha$ dla każdego $k \in \mathbb{N}$. Niech 
    \[
        c_k = k\alpha
    \]
    Wtedy
    \[
        \{ c_k: k \in \mathbb{N} \}
    \] nie ma elementu najmniejszego względem $\leq$,
    Wiemy, że $A\subseteq \mathbb{N}^n$, z wniosku z lematu Dicksona, $A$ ma skończenie wiele elemetów $\sqsubseteq$--minimalnych\footnote{$\sqsubseteq$ jest porządkiem produktowym.}.\\
    Niech $\{\alpha_1, \dots, \alpha_m\}$ to zbiór tych elementów. Numerujemy je tak, by $\alpha_1 < \alpha_2 < \dots < \alpha_m$. Pokażemy, że $\alpha_1$ jest elementem najmniejszym w $A$. \\
    Załóżmy, że $\beta < \alpha_1$ oraz $\beta \in A$. Ponieważ $\beta \notin M$, to $\exists i:\neg \alpha_i \sqsubseteq \beta$. Zatem
    \[
        \exists 0 \sqsubseteq \gamma: \beta + \gamma = \alpha_i
    \]
    co jest sprzeczne z założeniem, że $\alpha_i$ jest elementem minimalnym.\\
\end{proof}
\begin{definicja}[Leading Term, Monomial, Coefficient]~\\
    Niech $f \in K[\bar{x}]$, $f = \sum_{\alpha} a_\alpha \bar{x}^\alpha$. Najwyższym wyrazem $f$ względem porządku jednomianowego $\leq$ jest $LT(f) = a_\beta \bar{x}^\beta$, gdzie $\beta$ jest minimalnym elementem w zbiorze $\{\alpha: a_\alpha \neq 0\}$.
    \begin{itemize}
        \item $LT(f) = a_\beta \cdot \bar{x}^\beta$ -- leading term
        \item $LM(f) = \bar{x}^\beta$ -- leading monomial
        \item $LC(f) = a_\beta$ -- leading coefficient
    \end{itemize}
\end{definicja}
\begin{definicja}[Polynomial Quotient Reminder]~\\
    Niech $f, g_1, \dots, g_n \in K[\bar{x}]$, $\forall_{1\leq i \leq n}: g_i \neq 0$, $f$ jest zdefiniowane jako:
    \[
        f(\bar{x}) = \sum_{i=1}^n q_i(\bar{x})g_i(\bar{x}) + r(\bar{x})
    \]
    gdzie $q_i, r \in K[\bar{x}]$ oraz $r = 0$ lub $LT(r)$ nie jest podzielny przez żaden z $LT(g_i)$.
    \[
        \text{PQR}(f, g_1, \dots, g_n) = (q_1, \dots, q_n, r)
    \]
\end{definicja}
\begin{przyklad}~\\
    Niech $f = x^2y + xy^2 + 1$, $g_1 = xy - 1$, $g_2 = y^2 - 1$ oraz porządek leksykograficzny. Wtedy
    \[
        \text{PQR}(f, g_1, g_2) = (x+y, 0, 2)
    \]
\end{przyklad}

Jak to działa? Idea polega na tym, że dzielimy $f$ przez $g_1, g_2, \dots, g_n$ w sposób podobny do dzielenia wielomianów jednowymiarowych, ale z uwzględnieniem porządku jednomianowego. W każdej iteracji sprawdzamy, czy najwyższy wyraz $r$ jest podzielny przez którykolwiek z $LT(g_i)$. Jeśli tak, to dzielimy i aktualizujemy $q_i$ oraz $r$. Proces ten powtarzamy, aż $r$ będzie równy zero lub jego najwyższy wyraz nie będzie podzielny przez żaden z $LT(g_i)$.
Algorytm ten bedzie wyglądał następująco:
\begin{algorithm}
    \caption{Polynomial Reduce}
    \begin{algorithmic}[1]
        \Procedure{PolynomialReduce}{$f, g_1, \dots, g_n$} \Comment{$f, g_i \in K[\bar{x}]$, $g_i \neq 0$}
        \State $p \gets f$
        \State $r \gets 0$
        \State $(q_1, \dots, q_n) \gets (0, \dots, 0)$
        \State $i \gets 1$
        \While{$p \neq 0$}
            \If{$\exists k: LT(g_k) | LT(p)$}
                \State $q_k \gets q_k + \frac{LT(p)}{LT(g_k)}$
                \State $p \gets p - \frac{LT(p)}{LT(g_k)}g_k$
            \Else
                \State $r \gets r + LT(p)$
                \State $p \gets p - LT(p)$
            \EndIf
        \EndWhile
        \State \Return $(q_1, \dots, q_n, r)$
        \EndProcedure
    \end{algorithmic}
\end{algorithm}

\begin{fakt}
    \[
        \forall i: LM(g_i \cdot q_i) \leq LM(f)
    \]
\end{fakt}
\begin{przyklad}
    Niech $f = x^3 -x^2 y^2, g = x-y$ oraz dany jest porządek leksykograficzny $x<y$. Wtedy
    \begin{align*}
        x^3 - x^2 y^2 & \xrightarrow{-x^2(x-y)} -x^2 y^2 + x^2 y \\
        & \xrightarrow{-x y^2(x-y)} -x y^2 + x y^3 \\
        & \xrightarrow{-xy(x-y)} -xy^3 + xy^2\\
        & \xrightarrow{-y^3(x-y)} xy^2 + y^4 \\
        & \xrightarrow{-y^2(x-y)}- y^4 + y^3 \\
        & \xrightarrow{r = -y^4 + y^3} 0
    \end{align*}
    Można to zapisać jako:
    \[
        f= (x-y)(x^2 - x y^2 +xy - y^3- y^2) + (-y^4 + y^3)
    \]
\end{przyklad}

\begin{przyklad}
    Niech $f = xy^2 + y^3, g_1 = xy - 1, g_2 = y^2 - 1$ oraz porządek leksykograficzny $x < y$. Wtedy
    \begin{align*}
        xy^2 + y^3 & \xrightarrow{-y(xy - 1)} y^3 + y \\
                   & \xrightarrow{-y(y^2 - 1)} 2y \\
                   & \xrightarrow{r = 2y} 0
    \end{align*}
    Czyli innymi słowy
    \[
        f = (xy - 1)(-y) + (y^2 - 1)(-y) + 2y
    \]
\end{przyklad}

\begin{definicja}[Ideały jednomianowe]~\\
    $\mathcal{I}$ nazywamy ideałem jednomianowym, jeśli 
    \[
        \exists A \subseteq \mathbb{N}^n: \mathcal{I} = \langle \bar{x}^\alpha: \alpha \in A \rangle
    \]
\end{definicja}
Zastanówmy się, co jeżeli $f \in \mathcal{I}$
\[
    f \in \langle \bar{x}^\alpha: \alpha \in A \rangle
\]
to znaczy, że $f$ można przedstawić jako kombinację liniową wielomianów $\bar{x}^\alpha$ dla $\alpha \in A$. W takim przypadku, najwyższy wyraz $f$ będzie podzielny przez co najmniej jeden z $\bar{x}^\alpha$, ponieważ $LT(f)$ musi być podzielny przez $LT(\bar{x}^\alpha)$ dla co najmniej jednego $\alpha \in A$. Oznacza to, że $LT(f)$ jest podzielny przez co najmniej jeden z $\bar{x}^\alpha$ dla $\alpha \in A$.

\[
    f = \sum_{\alpha \in A} q_\alpha \bar{x}^\alpha
\]
gdzie $q_\alpha \in K[\bar{x}]$ oraz tylko skończenie wiele $q_\alpha$ jest różne od zera. Wtedy $LT(f)$ jest podzielny przez co najmniej jeden z $\bar{x}^\alpha$, ponieważ $LT(f)$ musi być podzielny przez $LT(q_\alpha \bar{x}^\alpha)$ dla co najmniej jednego $\alpha \in A$. Ponieważ $LT(q_\alpha \bar{x}^\alpha) = LT(q_\alpha) \cdot LT(\bar{x}^\alpha)$, to $LT(f)$ jest podzielny przez $LT(\bar{x}^\alpha)$ dla co najmniej jednego $\alpha \in A$.\\
Podsumowując:
\[
    f = \sum_{\alpha \in A} \sum_{\beta \in A_{q_\alpha}} a_{\alpha, \beta} \bar{x}^\beta \cdot \bar{x}^\alpha
\]
\begin{uwaga}~\\
    \[
        x^\alpha \sqsubseteq x^{\alpha +\beta} 
    \]
\end{uwaga}

\begin{fakt}~\\
    Jeśli $x^\gamma \in \mathcal{I}$, to istnie $\alpha \in A$ taka, że $x^\alpha \sqsubseteq x^\gamma$.
\end{fakt}

\begin{definicja}[]~\\
    Niech $\kappa$ -- ciało oraz $\mathcal{F} \subseteq \Kappa[\bar{x}]$. Określamy
    \begin{itemize}
        \item $LT(\mathcal{F}) = \{ LT(f): f \in \mathcal{F} \}$
        \item $LM(\mathcal{F}) = \{ LM(f): f \in \mathcal{F} \}$
    \end{itemize}
\end{definicja}

\begin{fakt}~\\
    $LT(\langle \mathcal{F} \rangle) = \langle LM(\mathcal{F}) \rangle$
\end{fakt}

\begin{definicja}[Baza Gröbnera]~\\
    Bazą Gröbnera ideału $\mathcal{I}$ ($\mathcal{I} \neq \{0\}$) nazywamy $G \subseteq \kappa[\bar{x}]$ takie, że $|G| < \aleph_0$, $LT(\langle G \rangle) = LT(\mathcal{I})$ oraz $G \subseteq \mathcal{I}$.
\end{definicja}

\begin{przyklad}~\\
    Niech $\mathcal{F} = \{x^3-2xy, x^2y -2y^2+x\}$ z porządkiem leksykograficznym $y<x$
    Rozważmy \[
        \mathcal{I} = \langle \mathcal{F} \rangle
    \]
    Wiemy, że \[
        LT(\mathcal{I}) = \langle x^3, x^2y \rangle
    \] 
    Zauważmy, że
    \begin{align*}
        y(x^3 - 2xy) - x(x^2y - 2y^2 + x) & = -2x^2y + 2xy^2 - x^2 \\
                                          & = -x^2 \in \mathcal{I}
    \end{align*}
    a więc $x^2 \in \mathcal{I}$, ale $x^2 \notin \langle \mathcal{F} \rangle$. Zatem $\mathcal{F}$ nie jest bazą Gröbnera dla $\mathcal{I}$.

\end{przyklad}

\begin{twierdzenie}~\\
    Każdy niepusty ideał $\mathcal{I} \subseteq \kappa[\bar{x}]$ ma bazę Gröbnera.
\end{twierdzenie}
\begin{proof}~\\
    Kożystając z wcześniejszego faktu, że $\langle LT(\mathcal{I}) \rangle  = \langle LM(\mathcal{I}) \rangle$ -- ideał jednomianowy. Ma on postać
    $\langle \{x^\alpha: \alpha \in A \subseteq \mathbb{N}^n\} \rangle$. Niech $A_0$ -- zbiór elementów minimalnych w $A$, zauważmy, że $A_0$ jest zbiorem skończonym z lemat Dicksona. Udowodnimy, że $\mathcal{A} = \{ x^\alpha : \alpha \in A_0\}$ jest bazą Gröbnera dla $\mathcal{I}$.\\
    Oczywiste jest, ze \[
    \mathcal{A} \subseteq \mathcal{I}\]
   Weźmy $f \in \mathcal{I}\setminus \{0\}$. Wtedy 
    \[
        LT(f) = c \cdot x^\xi \in \langle LM(\mathcal{I}) \rangle 
    \]
    Co z kolei oznacza, że
    \[
        x^\xi \in LM(\mathcal{I}) \equiv \xi \in A
    \]
    a więc
    \[
        \exists \beta: \exists \alpha \in A_0: \xi = \alpha + \beta
    \]
    czyli
    \[
        c x^\xi \in \langle x^\alpha \rangle \subseteq \langle \mathcal{A} \rangle
    \]
    Zatem $LT(f) \in \langle \mathcal{A} \rangle$. 
\end{proof}

\begin{twierdzenie}~\\
    Jeśli $G$ jest bazą Gröbnera dla $\mathcal{I}$, to $\langle G \rangle = \mathcal{I}$
\end{twierdzenie}
\begin{proof}~\\
    Weźmy $f \in \mathcal{I}$\setmius \{0\}, $|G| = s, G=\{g_i: 1\leq i \leq s\}$. Używając wcześniej stworzonego algorytmu \texttt{PolynomialReduce}, możemy przedstawić $f$ jako
    \[
        PolynomialReduce(f, G) \rightarrow (q_1, \dots, q_n, r)
    \]
    \[
        f = \underbrace{\sum_{i=1}^s q_i g_i}_{\in \langle G \rangle} + r
    \]
    Jeśli $r = 0$, to $f \in \langle G \rangle$.Jeśli $r \neq 0$, to żadne z termów $r$ nie dzieli się przez żadne z $LT(g_i) \quad (\bigstar)$.\\
    Skoro $G$ jest bazą Gröbnera dla $\mathcal{I}$, to $\langle LT(\mathcal{I}) \rangle = \langle LT(G) \rangle$, zatem 
    \[
        LT(r) = LT(f-\sum_{i=1}^s q_i g_i) \in \langle LT(\mathcal{I}) \rangle
    \]
    Stąd 
    \[
        LT(r) \in \langle LT(\mathcal{I}) \rangle = \langle LT(G) \rangle
    \]
    Ale z $(\bigstar)$ mamy \Rightarrow\!\Leftarrow
\end{proof}
\begin{wniosek}[Twierdzenie Hilberta]~\\
    Niepusty ideał ma skończoną bazę
\end{wniosek}
\begin{wniosek}[Własność rosnących łańcuchów (ACC)]~\\
    Załózmy, że mamy rosnący ciąg ideałów $\mathcal{I}_1 \subseteq \mathcal{I}_2 \subseteq \dots \subseteq \kappa[\bar{x}]$. Wtedy 
    \[
        \exists n \in \mathbb{N}: \mathcal{I}_n = \mathcal{I}_{n+1} = \dots
    \]
\end{wniosek}
\begin{definicja}~\\
    Pierścień $\mathcal{R}$ o własności ACC nazywamy pierścieniem noetherowskim.
\end{definicja}
\begin{przyklad}~\\
    $\kappa[x_1, x_2, \dots]$ jest pierścieniem nienoetherowskim
    \[
    \langle x_1 \rangle \nsubseteq \langle x_1, x_2 \rangle \nsubseteq \langle x_1, x_2, x_3 \rangle \nsubseteq \dots
\] można tak dokładać zmienne w nieskończoność.
\end{przyklad}

\end{document}
